\section{Methodology}
\label{sec:method}

Given a text prompt $y$, we generate an image by optimizing a latent
$\mathbf{x}_0\in\mathbb{R}^{1\times4\times64\times64}$ initialized from
$\mathcal{N}(\mathbf{0},\mathbf{I})$.
Stable Diffusion 2.1-base provides the pretrained noise predictor
$\hat{\boldsymbol{\epsilon}}_\theta(\mathbf{x}_t,y,t)$.
Its base parameters remain fixed, and the final latent is decoded into an RGB image using the
pretrained VAE.

Let $t$ denote a diffusion timestep and $\bar\alpha_t$ the cumulative product of the noise-schedule
coefficients.
Gaussian forward noising produces
\begin{equation}
 \mathbf{x}_t=\sqrt{\bar\alpha_t}\,\mathbf{x}_0
   +\sqrt{1-\bar\alpha_t}\,\boldsymbol{\epsilon}_t,
 \qquad \boldsymbol{\epsilon}_t\sim\mathcal{N}(\mathbf{0},\mathbf{I}).
 \label{eq:forward}
\end{equation}
We obtain the coefficients from the model's DDIM scheduler, with timestep bounds $t_{\min}=20$ and
$t_{\max}=980$ out of 1000 training timesteps.
The sampled noise $\boldsymbol{\epsilon}_t$ is distinguished from the predicted noise
$\hat{\boldsymbol{\epsilon}}_\theta$.

The implementation applies classifier-free guidance (CFG) by combining predictions conditioned on
the positive prompt $y$ and negative prompt $y_-$, with guidance weight $\omega$:
\begin{equation}
\begin{aligned}
 \hat{\boldsymbol{\epsilon}}_{\mathrm{CFG}}(\mathbf{x}_t,y,t)
 ={}&\hat{\boldsymbol{\epsilon}}_\theta(\mathbf{x}_t,y_-,t)\\
 &+\omega\bigl[\hat{\boldsymbol{\epsilon}}_\theta(\mathbf{x}_t,y,t)
 -\hat{\boldsymbol{\epsilon}}_\theta(\mathbf{x}_t,y_-,t)\bigr].
\end{aligned}
 \label{eq:cfg}
\end{equation}
Both text embeddings are computed without gradients, and the two predictions are evaluated in a
single batched U-Net call.
Below, $\hat{\boldsymbol{\epsilon}}_\theta$ denotes the prediction after CFG with the guidance
setting specified for each method.

%-------------------------------------------------------------------------
\subsection{Score Distillation Sampling (SDS)}
\label{sec:sds}

\noindent\textbf{Concept.}
SDS~\cite{poole2023dreamfusion} uses the difference between predicted and sampled noise to update a
parameterized representation.
Let $\boldsymbol{\theta}$ denote the optimized representation parameters, $c$ a camera viewpoint,
and $\mathbf{x}_0=g(\boldsymbol{\theta};c)$ the rendered image.
Omitting the noise predictor Jacobian gives
\begin{equation}
 \nabla_{\boldsymbol{\theta}}\mathcal{L}_{\mathrm{SDS}}(\boldsymbol{\theta})
 =\bigl(\hat{\boldsymbol{\epsilon}}_\theta(\mathbf{x}_t,y,t)
        -\boldsymbol{\epsilon}_t\bigr)
   \frac{\partial\mathbf{x}_t}{\partial\boldsymbol{\theta}}.
 \label{eq:sds_gradient}
\end{equation}
The subscript $\theta$ on the noise predictor identifies the pretrained network, while
$\boldsymbol{\theta}$ in the gradient denotes the parameters being optimized.

\noindent\textbf{Implementation.}
At each iteration, we sample an integer $t$ uniformly from $[20,980)$, draw Gaussian noise, and
construct $\mathbf{x}_t$ using \cref{eq:forward}.
Our implementation additionally weights the noise residual by $w(t)=1-\bar\alpha_t$ and applies the
weighted residual directly as the latent gradient.
The noise prediction runs under \texttt{torch.no\_grad()}.
We call \texttt{detach()} on the weighted residual and return its elementwise product with the
latent, summed to a scalar.
Backpropagation through this scalar supplies the chosen gradient without differentiating through the
U-Net.

\begin{lstlisting}
def get_sds_loss(self, latents, text_embeddings, guidance_scale=7.5):
    t = torch.randint(self.min_step, self.max_step, (latents.shape[0],), device=latents.device)
    noise = torch.randn_like(latents)
    alpha_bar_t = self.alphas[t].view(-1, 1, 1, 1)
    xt = torch.sqrt(alpha_bar_t) * latents + torch.sqrt(1 - alpha_bar_t) * noise
    with torch.no_grad():
        noise_pred = self.get_noise_preds(xt, t, text_embeddings, guidance_scale)
    w = 1 - alpha_bar_t
    grad = (w * (noise_pred - noise)).detach()
    loss = (grad * latents).sum()
    return loss
\end{lstlisting}

%-------------------------------------------------------------------------
\subsection{Variational Score Distillation (VSD)}
\label{sec:vsd}

\noindent\textbf{Concept.}
VSD~\cite{wang2023prolificdreamer} learns a noise predictor for the current generated samples to
address over-smoothing and over-saturation.
The sampled-noise term in SDS is replaced by the LoRA-adapted prediction:
\begin{equation}
\begin{aligned}
 \nabla_{\boldsymbol{\theta}}\mathcal{L}_{\mathrm{VSD}}(\boldsymbol{\theta})
 ={}&\bigl(\hat{\boldsymbol{\epsilon}}_\theta(\mathbf{x}_t,y,t)\\
     &-\hat{\boldsymbol{\epsilon}}_{\mathrm{LoRA}}(\mathbf{x}_t,y,t)\bigr)
       \frac{\partial\mathbf{x}_t}{\partial\boldsymbol{\theta}}.
\end{aligned}
 \label{eq:vsd_gradient}
\end{equation}
We optimize one latent per prompt.
LoRA is trained with a denoising mean squared error between its predicted noise and the Gaussian
noise used to construct $\mathbf{x}_t$.

\noindent\textbf{Implementation.}
We disable the adapters for the frozen teacher prediction, then enable them for the trainable
prediction on the same noisy latent and timestep.
The adapted predictor uses $\omega=1$ so that CFG reduces to the positive-condition prediction,
allowing LoRA to learn the conditional denoising prediction without guidance amplification.

As in SDS, the latent update uses the noise residual weighted by $1-\bar\alpha_t$.
The residual is detached before forming the scalar used for latent optimization.
The adapter's denoising loss is added with \texttt{lora\_loss\_weight=1.0}, and its input latent is
detached.
These operations ensure that the distillation term updates the latent and the denoising term updates
only LoRA.

\begin{lstlisting}
def get_vsd_loss(self, latents, text_embeddings, guidance_scale=7.5, lora_loss_weight=1.0):
    noise = torch.randn_like(latents)
    t = torch.randint(self.min_step, self.max_step, (latents.shape[0],), device=latents.device)
    alpha_bar_t = self.alphas[t].view(-1, 1, 1, 1)
    xt = torch.sqrt(alpha_bar_t) * latents + torch.sqrt(1 - alpha_bar_t) * noise
    with torch.no_grad():
        self.unet.disable_adapters()
        noise_pred = self.get_noise_preds(xt, t, text_embeddings, guidance_scale)
        self.unet.enable_adapters()
    noise_lora = self.get_noise_preds(xt.detach(), t, text_embeddings, 1.0)
    loss_lora = F.mse_loss(noise_lora, noise)
    w = 1 - alpha_bar_t
    grad = (w * (noise_pred - noise_lora)).detach()
    loss = (grad * latents).sum() + loss_lora * lora_loss_weight
    return loss
\end{lstlisting}

%-------------------------------------------------------------------------
\subsection{Score Distillation via Inversion (SDI)}
\label{sec:sdi}

\noindent\textbf{Concept.}
SDI~\cite{lukoianov2024sdi} obtains a noisy sample by DDIM inversion~\cite{song2021ddim} instead of
injecting independently sampled Gaussian noise.
The current latent is updated by minimizing the mean squared error to a denoised target:
\begin{equation}
 \mathcal{L}_{\mathrm{SDI}}
 =\operatorname{mean}\!\left[(\mathbf{x}_0-\mathbf{x}_{0|t})^2\right].
 \label{eq:sdi_reconstruction_loss}
\end{equation}
Here, $\mathbf{x}_{0|t}$ is the cached clean estimate reconstructed directly from the inverted
$\mathbf{x}_t$ at the same timestep $t$.
The square is elementwise, and the mean averages over all latent elements.
The timestep follows a rounded linear schedule from 980 to 20 over the optimization iterations
rather than being sampled independently.

\noindent\textbf{DDIM inversion.}
We implement the inversion loop explicitly without calling a library inversion utility.
The loop starts from the current clean latent and traverses rounded, evenly spaced timesteps from
zero to the target timestep.
At each step, it predicts noise, estimates the clean latent, and combines that estimate with the
predicted noise using the coefficients of the next scheduled timestep.
The initial coefficient is set to 1 because the input is clean.
Inversion is deterministic: the accepted \texttt{eta} argument is unused, as permitted by the
assignment.

\begin{lstlisting}
@torch.no_grad()
def invert_noise(self, latents, target_t, text_embeddings, guidance_scale=-7.5, n_steps=10, eta=0.3):
    if n_steps < 1:
        raise ValueError("n_steps must be at least 1")
    target_t = torch.as_tensor(target_t, device=latents.device).expand(latents.shape[0])
    stepsize = target_t / n_steps
    timesteps = (
        torch.arange(n_steps + 1, device=latents.device)[:, None] * stepsize
    ).round().long()
    alphas = self.alphas[timesteps].to(latents.dtype).view(
        n_steps + 1, latents.shape[0], 1, 1, 1
    )
    alphas[0] = 1
    sqrt_alphas = alphas.sqrt()
    sqrt_one_minus_alphas = (1 - alphas).sqrt()
    for i in range(n_steps):
        noise_pred = self.get_noise_preds(
            latents, timesteps[i], text_embeddings, guidance_scale
        )
        pred_x0 = (
            latents - sqrt_one_minus_alphas[i] * noise_pred
        ) / sqrt_alphas[i]
        latents = (
            sqrt_alphas[i + 1] * pred_x0
            + sqrt_one_minus_alphas[i + 1] * noise_pred
        )
    return latents
\end{lstlisting}

\noindent\textbf{SDI Implementation.}
The implemented target is the clean estimate reconstructed directly from the inverted noisy latent:
\begin{equation}
 \mathbf{x}_{0|t}
 =\frac{\mathbf{x}_t-\sqrt{1-\bar\alpha_t}\,
         \hat{\boldsymbol{\epsilon}}_\theta(\mathbf{x}_t,y,t)}
        {\sqrt{\bar\alpha_t}}.
 \label{eq:sdi_target}
\end{equation}

\begin{lstlisting}
def get_sdi_loss(
    self,
    latents,
    text_embeddings,
    guidance_scale=7.5,
    current_iter=0,
    total_iters=500,
    inversion_guidance_scale=-7.5,
    inversion_n_steps=10,
    inversion_eta=0.3,
    update_interval=25,
):
    B = latents.shape[0]
    progress = min(max(current_iter / max(total_iters - 1, 1), 0.0), 1.0)
    timestep = round(self.max_step + (self.min_step - self.max_step) * progress)
    t = torch.full((B,), timestep, device=latents.device, dtype=torch.long)
    should_update = (current_iter % update_interval == 0) or not hasattr(self, 'sdi_target')
    if should_update:
        with torch.no_grad():
            latents_noisy = self.invert_noise(
                latents, t, text_embeddings,
                guidance_scale=inversion_guidance_scale,
                n_steps=inversion_n_steps,
                eta=inversion_eta
            )
            noise_pred = self.get_noise_preds(latents_noisy, t, text_embeddings, guidance_scale)
            alpha_bar_t = self.alphas[t].to(latents.dtype).view(-1, 1, 1, 1)
            target = (
                latents_noisy - torch.sqrt(1 - alpha_bar_t) * noise_pred
            ) / torch.sqrt(alpha_bar_t)
            self.sdi_target = target.detach()
    loss = F.mse_loss(self.sdi_target, latents)
    return loss
\end{lstlisting}
